% GENERATED FILE. Edit canonical lessons, then run npm run generate:books.
% canonical-source-hash: fnv1a64:852fcbbdcbf10b6f
% canonical-lessons: SA-C67-sleep, SA-C67-hunger, SA-C67-thirst, SA-C67-rest, SA-C67-ease

\chapter{Tired, Hungry, Thirsty}
\label{ch:sa-tired-hungry-thirsty}
\hlchaptermodality{67}

\begin{chapteropening}
\textbf{By the end of this chapter:} \emph{I can say I am sleepy, hungry or thirsty, ask for rest, and say I am at ease.}

Complete the last lesson of chapter 67: i can say I am sleepy, hungry or thirsty, ask for rest, and say I am at ease.
\end{chapteropening}

\section[nidrā]{\sk{निद्रा} (nidrā) — sleep}

\label{lesson:SA-C67-sleep}



\begin{quote}
Before the new one: say the Sanskrit for an illness, then the Sanskrit for health.
\end{quote}

\subsection*{You'll want to know: \sk{निद्रा}}
\textbf{\sk{निद्रा}} (\emph{nidrā}) — \textquotedblleft{}sleep\textquotedblright{}. It comes at \textbf{\sk{रात्रिः}}.

\begin{grammarlens}[title={What you've built}]
What the night is for.
\end{grammarlens}

\subsection*{Your turn}
\begin{itemize}
\raggedright
  \item \emph{Say it:} \emph{nidrā}
  \item \emph{Say it:} \emph{nidrā}, once more
  \item \emph{Say it:} say \emph{nidrā}
  \item \emph{Recall:} say the Sanskrit for an illness, then the Sanskrit for health, then say \emph{nidrā} again
\end{itemize}

\begin{tcolorbox}[breakable,colback=teal!4,colframe=teal!35!black,title={Before you move on}]
What does \sk{निद्रा} mean? (\textquotedblleft{}Sleep\textquotedblright{}.) Say it once more.
\end{tcolorbox}
\section[kṣudhā]{\sk{क्षुधा} (kṣudhā) — hunger}

\label{lesson:SA-C67-hunger}



\begin{quote}
Before the new one: say the Sanskrit for health, then the Sanskrit for sleep.
\end{quote}

\subsection*{You'll want to know: \sk{क्षुधा}}
\textbf{\sk{क्षुधा}} (\emph{kṣudhā}) — \textquotedblleft{}hunger\textquotedblright{}. \textbf{\sk{मम} \sk{क्षुधा} \sk{अस्ति}} — \textquotedblleft{}I am hungry\textquotedblright{}.

\begin{grammarlens}[title={What you've built}]
The need that food answers.
\end{grammarlens}

\subsection*{Your turn}
\begin{itemize}
\raggedright
  \item \emph{Say it:} \emph{kṣudhā}
  \item \emph{Say it:} \emph{kṣudhā}, once more
  \item \emph{Say it:} say \emph{mama kṣudhā asti}
  \item \emph{Recall:} say the Sanskrit for health, then the Sanskrit for sleep, then say \emph{kṣudhā} again
\end{itemize}

\begin{tcolorbox}[breakable,colback=teal!4,colframe=teal!35!black,title={Before you move on}]
What does \sk{क्षुधा} mean? (\textquotedblleft{}Hunger\textquotedblright{}.) Say it once more.
\end{tcolorbox}
\section[pipāsā]{\sk{पिपासा} (pipāsā) — thirst}

\label{lesson:SA-C67-thirst}



\begin{quote}
Before the new one: say the Sanskrit for sleep, then the Sanskrit for hunger.
\end{quote}

\subsection*{You'll want to know: \sk{पिपासा}}
\textbf{\sk{पिपासा}} (\emph{pipāsā}) — \textquotedblleft{}thirst\textquotedblright{}. \textbf{\sk{पानीयम्}} answers it.

\begin{grammarlens}[title={What you've built}]
The need that water answers.
\end{grammarlens}

\subsection*{Your turn}
\begin{itemize}
\raggedright
  \item \emph{Say it:} \emph{pipāsā}
  \item \emph{Say it:} \emph{pipāsā}, once more
  \item \emph{Say it:} say \emph{mama pipāsā asti}
  \item \emph{Recall:} say the Sanskrit for sleep, then the Sanskrit for hunger, then say \emph{pipāsā} again
\end{itemize}

\begin{tcolorbox}[breakable,colback=teal!4,colframe=teal!35!black,title={Before you move on}]
What does \sk{पिपासा} mean? (\textquotedblleft{}Thirst\textquotedblright{}.) Say it once more.
\end{tcolorbox}
\section[viśrāmaḥ]{\sk{विश्रामः} (viśrāmaḥ) — rest}

\label{lesson:SA-C67-rest}



\begin{quote}
Before the new one: say the Sanskrit for hunger, then the Sanskrit for thirst.
\end{quote}

\subsection*{You'll want to know: \sk{विश्रामः}}
\textbf{\sk{विश्रामः}} (\emph{viśrāmaḥ}) — \textquotedblleft{}rest\textquotedblright{}, a pause from work or from a \textbf{\sk{यात्रा}}.

\begin{grammarlens}[title={What you've built}]
What a tired traveller wants.
\end{grammarlens}

\subsection*{Your turn}
\begin{itemize}
\raggedright
  \item \emph{Say it:} \emph{viśrāmaḥ}
  \item \emph{Say it:} \emph{viśrāmaḥ}, once more
  \item \emph{Say it:} say \emph{viśrāmaḥ}
  \item \emph{Recall:} say the Sanskrit for hunger, then the Sanskrit for thirst, then say \emph{viśrāmaḥ} again
\end{itemize}

\begin{tcolorbox}[breakable,colback=teal!4,colframe=teal!35!black,title={Before you move on}]
What does \sk{विश्रामः} mean? (\textquotedblleft{}Rest\textquotedblright{}.) Say it once more.
\end{tcolorbox}
\section[sukham]{\sk{सुखम्} (sukham) — ease, happiness}

\label{lesson:SA-C67-ease}



\begin{quote}
Before the new one: say the Sanskrit for thirst, then the Sanskrit for rest.
\end{quote}

\subsection*{You'll want to know: \sk{सुखम्}}
\textbf{\sk{सुखम्}} (\emph{sukham}) — \textquotedblleft{}ease\textquotedblright{}, \textquotedblleft{}happiness\textquotedblright{}. \textbf{\sk{सुखम्} \sk{अस्ति}} — \textquotedblleft{}all is well\textquotedblright{}.

\begin{grammarlens}[title={What you've built}]
How things are after rest.
\end{grammarlens}

\subsection*{Your turn}
\begin{itemize}
\raggedright
  \item \emph{Say it:} \emph{sukham}
  \item \emph{Say it:} \emph{sukham}, once more
  \item \emph{Say it:} say \emph{sukham asti}
  \item \emph{Recall:} say the Sanskrit for thirst, then the Sanskrit for rest, then say \emph{sukham} again
\end{itemize}

\begin{tcolorbox}[breakable,colback=teal!4,colframe=teal!35!black,title={Before you move on}]
What does \sk{सुखम्} mean? (\textquotedblleft{}Ease, happiness\textquotedblright{}.) Say it once more.
\end{tcolorbox}
